\section{Introducción}

Las desigualdades territoriales representan una de las expresiones más evidentes
del desarrollo económico desigual. Aun cuando las regiones de un mismo país
comparten un mismo marco institucional y macroeconómico, pueden seguir
trayectorias muy distintas en términos de ingreso, productividad, educación,
salud y condiciones de vida. Esto significa que el crecimiento del producto por
habitante no siempre se traduce en mejoras similares para toda la población.
En este sentido, el análisis del desarrollo regional requiere considerar tanto los
aspectos económicos como los factores sociales que influyen en las oportunidades
y capacidades de las personas.

La teoría neoclásica del crecimiento plantea que las economías con menores
niveles iniciales de capital e ingreso pueden crecer a un ritmo más acelerado,
debido a los rendimientos marginales decrecientes del capital
\parencite{solow1956}. Bajo este enfoque, la convergencia beta ocurre cuando las
regiones inicialmente rezagadas presentan tasas de crecimiento superiores a las
de aquellas con mayores niveles de desarrollo. Esta idea fue contrastada de
manera empírica por \textcite{barro_sala1991} y
\textcite{barro_sala1992}, y posteriormente ampliada con la incorporación del
capital humano y otros factores estructurales
\parencite{mankiw1992}.

Un coeficiente negativo entre el nivel inicial y el crecimiento posterior no
demuestra, por sí mismo, que las brechas regionales estén disminuyendo.
\textcite{quah1993} advierte que las regresiones tradicionales de convergencia
pueden ocultar transformaciones en la distribución, persistencia de
desigualdades e incluso la formación de grupos de regiones con trayectorias
diferenciadas. De ahí la necesidad de complementar la estimación beta con
medidas de dispersión, densidades kernel y otros procedimientos que permitan
examinar la evolución de la distribución en su conjunto.

La dimensión espacial añade un elemento central al análisis. Las regiones no
siguen trayectorias independientes, pues mantienen relaciones económicas,
productivas, laborales y geográficas con los territorios próximos. La
localización de la infraestructura, la concentración de actividades
económicas, el acceso a los mercados y las economías de aglomeración pueden
reforzar las ventajas de determinados espacios \parencite{krugman1991}.
Prescindir de estas interacciones podría generar estimaciones incompletas y
lecturas poco representativas de la organización territorial de la economía
\parencite{anselin1988,rey_montouri1999}.

La interpretación de las desigualdades regionales también cambia según el
indicador seleccionado. El ingreso per cápita aproxima la capacidad económica
promedio de un territorio, aunque no resume adecuadamente sus condiciones de
salud, educación y bienestar. Desde el enfoque de las capacidades, el
desarrollo se relaciona con la ampliación de las libertades reales de las
personas y no únicamente con el incremento de sus ingresos
\parencite{sen1999}. Bajo esta perspectiva, la investigación examina
conjuntamente el PBI real per cápita relativo y el Índice de Desarrollo Humano.
La comparación permite establecer si el acercamiento económico entre
departamentos estuvo acompañado por una reducción semejante de las brechas
sociales.

La evidencia internacional indica, además, que la convergencia puede presentar
intensidades distintas dentro de un mismo país. El estudio de
\textcite{duran2024} sobre las regiones de Turquía combina indicadores
socioeconómicos, econometría espacial y regresiones geográficamente ponderadas,
y encuentra que los parámetros estimados varían territorialmente. El presente
trabajo retoma esa estrategia para el caso peruano, con adecuaciones derivadas
de la disponibilidad de información, el periodo de estudio y la escala
departamental de las variables.

En el Perú, las diferencias acumuladas en infraestructura, capital humano,
estructura productiva, urbanización y acceso a los mercados justifican una
lectura explícitamente territorial. Algunos departamentos concentran
actividades de mayor productividad y una oferta más amplia de servicios,
mientras otros todavía enfrentan restricciones estructurales que limitan la
generación de ingresos y la ampliación de capacidades. Una evaluación basada
solo en promedios nacionales podría ocultar tanto la magnitud de estas brechas
como la diversidad de las trayectorias departamentales.

Sobre esta base, la investigación plantea la siguiente pregunta: ¿existió
convergencia regional en desarrollo humano y PBI real per cápita relativo
entre los departamentos del Perú durante 2007--2019, y en qué medida dicho
proceso presentó heterogeneidad espacial? El objetivo general es analizar las
disparidades socioeconómicas y la convergencia regional a partir de la
distribución de los indicadores, la autocorrelación espacial y la variación
territorial de los coeficientes estimados.

La estrategia empírica combina estadísticos descriptivos, densidades kernel,
mapas temáticos, el índice global de Moran, regresiones de convergencia beta
mediante mínimos cuadrados ordinarios, modelos de error espacial,
especificaciones Durbin de error y regresiones geográficamente ponderadas. Los
hallazgos revelan trayectorias diferentes entre las dos dimensiones
consideradas. El PBI real per cápita relativo presenta evidencia de
convergencia beta, mientras que el IDH mantiene una relación positiva entre su
nivel inicial y el cambio posterior. La heterogeneidad territorial también es
más marcada en la convergencia económica que en la evolución del desarrollo
humano.

El artículo se divide en seis secciones. Luego de esta introducción, la segunda
sección desarrolla la revisión de literatura y los fundamentos teóricos. La
tercera presenta las fuentes de información y el procedimiento metodológico.
La cuarta reúne los resultados descriptivos, espaciales y econométricos. La
quinta discute los principales hallazgos, sus implicancias y las limitaciones
del estudio. La última sección expone las conclusiones.
